\clearpage\begingroup\NoSourceHead
\renewcommand{\thefootnote}{\arabic{footnote})}\setcounter{footnote}{13}
\ArithmeticHead{10}{L}{BOLZA, GAUSS UND DIE VARIATIONSRECHNUNG.}
\noindent\hspace*{\parindent}5. Die Methode der Variation der Grenzen, wie sie von \textsc{Ohm} und \textsc{Sarrus}
ausgebildet worden ist, lässt an Klarheit und Einfachheit nichts zu wünschen
\par\clearpage\SourceLeaf{468}
\ArithmeticHead{11}{R}{GESCHICHTE DER EXTREMA VON DOPPELINTEGRALEN VOR GAUSS.}
\noindent übrig. Sie leidet aber an einer Schwäche, auf die wir schon oben angespielt
haben (Nr. 2), und die zu Tage tritt, sobald man die Methode auf geometrische
Aufgaben anwenden will. Man sehe sich z. B. die geradezu ungeheuerlichen
Formeln an, die \textsc{Sarrus}\footnote{\textsc{Sarrus}, loc. cit. 10), S. 102—114.} nötig hat, um das einfache Problem des Körpers
von grösstem Volumen bei gegebener Oberfläche zu behandeln.

Der Grund dieser Schwäche scheint mir darin zu liegen, dass die Aufgabe,
die durch die \textsc{Ohm-Sarrus}sche Methode gelöst wird, im allgemeinen für
die Geometrie kein Interesse hat, während umgekehrt die Aufgaben, die für
die Geometrie Interesse haben, nicht unter die \textsc{Ohm-Sarrus}sche Theorie fallen.
Denn der Integrationsbereich des Integrals (8) ist der durch die Geraden $AB$:
$x=x_0$, $CD$: $x=x_1$ einerseits und die beiden Kurven $AC$: $y=y_0(x)$, $BD$:
$y=y_1(x)$ andererseits eingeschlossene Bereich der $x,y$-Ebene, wobei $x_0<x_1$
und $y_0(x)\EqualAboveLess y_1(x)$ für $x_0\EqualAboveLess x\EqualAboveLess x_1$ vorausgesetzt wird. Einen solchen Bereich
wollen wir der Kürze halber einen »Elementarbereich« nennen.

Dann ist auch für das variierte Integral $\bar J$ der Integrationsbereich ein
Elementarbereich $\bar A\bar C\bar D\bar B$ und die Aufgabe, mit deren Lösung sich die Methode
der Variation der Grenzen beschäftigt, lautet in geometrischer Fassung
(für den Fall des Doppelintegrals):

Unter allen Flächenstücken: $z=z(x,y)$, deren Projektionen auf die $x,y$-Ebene
Elementarbereiche sind, dasjenige zu bestimmen, das ein Extremum für
das Doppelintegral $J$ liefert.

Diese Beschränkung auf Elementarbereiche erscheint aber vom geometrischen
Standpunkt aus durchaus unnatürlich und gekünstelt, und wird wohl
nur in den seltensten Fällen einmal bei einer aus der Geometrie selbst entsprungenen
Aufgabe auftauchen.

Man kann nun versuchen, diesem Mangel dadurch abzuhelfen, dass man
dem Elementarbereich die Bedingung auferlegt, dass die Punkte $A$ und $B$, $C$
und $D$ zusammenfallen sollen, und dass überdies die beiden Kurven $y=y_0(x)$
und $y=y_1(x)$ die Ordinaten $x=x_0$ und $x=x_1$ in den Punkten $A=B$, bezw.
$C=D$ berühren, und analog für die variierten Kurven.

Dann tritt aber der Übelstand ein, dass die Ableitungen $dy_0/dx$ und $dy_1/dx$
für die Werte $x=x_0$ und $x=x_1$ unendlich werden. Das macht sich zwar
\par\smallskip\noindent\hfill{\small 2\textsuperscript{*}}\par
\clearpage\SourceLeaf{469}
\ArithmeticHead{12}{L}{BOLZA, GAUSS UND DIE VARIATIONSRECHNUNG.}
\noindent bei der Ableitung der Formel (9) noch nicht geltend, macht aber die Umformung
der ersten Variation durch partielle Integration, — wenigstens in
der Form, wie sie von \textsc{Ohm}, \textsc{Sarrus} und \textsc{Delaunay} gehandhabt worden
ist —, hinfällig. Das ist ein Einwand, der auch die sonst so vortrefflichen
Lehrbücher von \textsc{Jellett} und \textsc{Lindelöf-Moigno} unbefriedigend erscheinen lässt,
soweit es sich um geometrische Anwendungen der Variationsrechnung handelt.

Die Lösung der hier vorliegenden Schwierigkeit wird einerseits durch
Einführung des Begriffs des Linienintegrals, andererseits durch die Methode
der Variation der unabhängigen Variabeln geliefert, zu deren Betrachtung
wir nunmehr übergehen.

6. Die Methode der \textls[100]{Variation der unabhängigen Variabeln} knüpft
unmittelbar an die Definition des $\delta$-Algorithmus an, die \textsc{Lagrange}\footnote{\textsc{Lagrange}, loc. cit. 1), S. 336. Es verdient jedoch hervorgehoben zu werden, dass \textsc{Lagrange} in seinen ersten \textls[100]{brieflichen} Mitteilungen über seine Entdeckung an \textsc{Euler} vom Jahre 1755 die Variable $x$ noch nicht variiert, vgl. \textsc{Lagrange}, loc. cit. 6), S. 140, 148.} in seiner
ersten Veröffentlichung (1760/1) über Variationsrechnung zunächst für einfache
Integrale gegeben hat:

Unter $Z$ wird eine beliebige Funktion der Variabeln $x,y,z$ und einer
Anzahl ihrer \textls[100]{Differentiale} verstanden und dann $\delta Z$ definiert als das Differential
dieser Funktion $Z$ in Bezug auf die als unabhängige Variable betrachteten
Grössen $x,y,z,dx,dy,dz,d^2x,d^2y,d^2z,\ldots$ Dazu kommen dann
die Gleichungen
\[
\begin{gathered}
\delta\!\int Z=\int\delta Z,\\
\delta dx=d\delta x,\quad\delta d^2x=d^2\delta x,\quad\ldots
\end{gathered}
\]

Es unterliegt keinem Zweifel, dass \textsc{Lagrange} sich dabei die Variabeln
$x,y,z$ als Funktionen einer unbestimmt gelassenen unabhängigen Hilfsvariabeln,
die wir $t$ nennen wollen, gedacht hat, obgleich er dies nirgends ausdrücklich
ausgesprochen hat. Daher kann man bei \textsc{Lagrange} streng genommen noch
nicht von einer Variation der \textls[100]{unabhängigen} Variabeln sprechen, insofern
bei ihm $x,y,z$ gleichberechtigt auftreten und die wahre unabhängige Variable
$t$ nicht variiert wird.

Etwas anders liegen die Dinge bereits bei \textsc{Euler}\footnote{\textsc{Euler}, loc. cit. 2), Art. 1—174.}, der in den \textit{Institutiones
calculi integralis} (1770) eine ausführliche und systematische Darstellung und
\par\clearpage\SourceLeaf{470}
\ArithmeticHead{13}{R}{GESCHICHTE DER EXTREMA VON DOPPELINTEGRALEN VOR GAUSS.}
\noindent Weiterführung der \textsc{Lagrange}schen Variationsmethode gegeben hat. Da er im
Gegensatz zu \textsc{Lagrange} das zu variierende Integral im einfachsten Fall in
der Form
\[
\begin{gathered}
U=\int V(x,y,p,q,\ldots)\,dx,\\
p=\frac{dy}{dx},\quad q=\frac{dp}{dx},\quad\ldots
\end{gathered}
\]
annimmt, so erscheint die Variable $x$ so ausgesprochen in der Rolle der unabhängigen
Variabeln, dass man bei \textsc{Euler} füglich von einer Variation der
\textls[100]{unabhängigen} Variabeln reden kann, wenn er selbst sich auch $x$ und $y$
als Funktionen einer Hilfsvariabeln gedacht haben mag, wofür mancherlei
Anzeichen sprechen.

Nun hat zwar einerseits \textsc{Euler} durch geometrische Deutung des Variationsprozesses
dem \textsc{Lagrange}schen Algorithmus einen mehr konkreten Inhalt zu
geben gesucht, andererseits hat er aber doch auch zu einer rein mechanischen
Anwendung des $\delta$-Prozesses beigetragen; so, wenn er z. B. ohne jede weitere
Begründung schreibt
\[
\delta U=\int\delta(V\,dx)=\int(\delta V\,dx+V\delta dx),
\]
oder wenn er die Gleichung
\[
\delta\left(\frac{dy}{dx}\right)=\frac{dx\,\delta dy-dy\,\delta dx}{dx^2}
=\frac{d\delta y}{dx}-\frac{dy}{dx}\frac{d\delta x}{dx}
\]
mit den Worten begründet\footnote{\textsc{Euler}, loc. cit. 2), Art. 45.}: »Variatio quaesita $\delta p$ per notas differentiationis
regulas reperitur, dummodo loco signi differentiationis $\partial$ scribatur signum variationis
$\delta$«. Ganz unverhüllt ist dann später diese wörtliche, rein mechanische
Anwendung der \textsc{Lagrange}schen Definition des $\delta$-Algorithmus, — für die wir
in der Folge die Bezeichnung »\textit{formale Variation}« gebrauchen werden — bei
der die Bedeutung der Operation vollständig hinter dem Algorithmus zurückgetreten
ist, und bei der keinerlei Bezugnahme auf eine Hilfsvariable mehr
erkennbar ist, bei \textsc{Lacroix}\footnote{\textit{Traité du calcul différentiel et du calcul intégral}, Bd. II, 2. Aufl., Paris 1814. Vgl. insbesondere Art. 845, 861, 862.} aufgetreten und hat von da aus den unheilvollsten
Einfluss auf die Weiterentwicklung der Variationsrechnung ausgeübt. Denn
diese wörtliche Anwendung der \textsc{Lagrange}schen Definition des $\delta$-Algorithmus
\par\clearpage\SourceLeaf{471}
\ArithmeticHead{14}{L}{BOLZA, GAUSS UND DIE VARIATIONSRECHNUNG.}
\noindent bei gleichzeitiger Variation der unabhängigen Variabeln ist \textls[100]{falsch}, wie schon
die direkte Anwendung auf die zweite Ableitung $d^2y/dx^2$ zeigt, und sie musste,
konsequent durchgeführt, früher oder später zu falschen Resultaten führen,
wie dies denn auch bei der Variation der Doppelintegrale tatsächlich eingetreten
ist.
\par\clearpage\endgroup
